\scanpage{9}
\mspage{PU 79}

sans voix ; \deleted{\unsure{et}} sa raison d'être (semblerait-il) n'est pas d'être \uline{entendu}, mais d' \mbox{ê\,t\,r\,e}. Il est le \mbox{C\,r\,é\,a\,t\,e\,u\,r}, celui qui crée sans témoins, avant qu'aucune créature encore n'entre en dialogue avec lui \handwritten{ou} \deleted{même} se réjouisse avec lui de ses oeuvres.

On dit que la langue des images est la langue de l'enfant\handwritten{,} \deleted{lui décrivant}\handwritten{\inserted{La ``langue'' par laquelle il appréhende}} le monde qui l'entoure. J'ai \deleted{\handwritten{\unsure{bien}}} oublié mon enfance, mais quelque chose pourtant me dit que c'est \deleted{bien} vrai, que c'est \handwritten{\inserted{bien là aussi}} la langue de ma propre enfance. Je ne sais si je la retrouverai un jour. Quelqu'un en moi pourtant la parle, cette langue; comme je la parlais, spontanément et sans effort - avant \deleted{de}\handwritten{\inserted{que je}} l'enterre\deleted{r}, un jour, et l'oublie\deleted{r}. Quelqu'un en moi la parle, mais c'est rare que je prenne le loisir d'écouter.

Je ne suis pas le seul, à l'écouter rarement, voire, jamais. Nous avons si bien appris à ne pas l'écouter, et nous sommes si bien embarqués sur le bateau ``Pensée'' à majuscules, alias ``Abstraction'', à la coque étanche de \mbox{m\,o\,t\,s} finement assemblés et rivés \handwritten{!} \deleted{\unsure{même|alors}} \handwritten{A}lors même que nous voudrions trouver le chemin d'un retour vers la langue oubliée de notre enfance, vers la source des rires et des pleurs, et vers les détresses oubliées, peut-être, que \handwritten{j}amais, jamais des mo\handwritten{t}s ne pourront dire - \deleted{\unsure{que}} ce chemin semble perdu à jamais \ldots

% archive sheet number: 8
